\section{Source-to-Event Expansion}\label{sec:compression}

A major objective of Motivo is to provide a highly expressive interface for music composition, reducing the
manual effort
required to construct complex arrangements.
We evaluate this expressiveness by analyzing the ``compression ratio'', the difference between the amount of code
written by the composer and the number of low-level MIDI events generated by the compiler.

\subsection{Algorithmic Conciseness}\label{detail-subsec:algorithmic-conciseness}

Traditional MIDI sequencers and digital audio workstations require the composer to manually place every note on a grid.
In contrast, Motivo allows composers to define structural rules and relationships.

Consider the following drum pattern from a representative studio example:
\begin{lstlisting}[language=Motivo, caption={A generative drum pattern using control flow.},label={lst:evaluation-drum-pattern}]
track percussion using drums {
    pattern groove() {
        for (int i = 0; i < 16; i = i + 1) {
            play hihat from i;
            if (i % 8 == 0) { play kick from i; }
            if (i % 8 == 4) { play snare from i; }
        }
    }
    loop (3) { play groove(); }
}
\end{lstlisting}

This block encapsulates a complex rhythm.
The \texttt{for} loop defines the underlying hi-hat grid, while the modulo-based \texttt{if} conditions procedurally
generate the kick and snare syncopation.
The \texttt{loop (3)} statement then instantiates this 16-beat pattern three times across the timeline.

\subsection{The Compression Ratio}\label{detail-subsec:the-compression-ratio}

To achieve the same result in a traditional MIDI editor, the composer would have to manually draw 48 hi-hat notes, 6
kick drum notes, and 6 snare drum notes (60 individual note events in total), which translates to 120 low-level MIDI
messages (\texttt{Note-On} and \texttt{Note-Off}).

Motivo abstracts these 120 MIDI messages into just a handful of declarative rules.
The resulting ``compression ratio'' is substantial on this example.
As the composition grows (for example, looping the groove 32 times instead of 3), the size of Motivo source code
remains completely unchanged, while the output scales to hundreds or thousands of events.

This expansion factor varies enormously with how repetitive a program is, and the benchmark suite spans both extremes.
The synthetic fixtures are written to maximize repetition: the flat scaling loop reaches about $2{,}200$ events per
source line at the 20{,}000-event cap (Table~\ref{tab:scaling-stages}), and the structural stress fixtures, each
amplifying a single construct, expand comparably.
Real compositions sit at the other end: the four pieces in Section~\ref{sec:general-case} expand only $1.9$ to $5.1$
events per source line (Table~\ref{tab:general-case}), because genuine music interleaves largely unique material rather
than repeating one statement.
The leverage of the language therefore depends on the piece: it removes manual event placement in every case, but the
dramatic ratios appear only where real structural repetition exists.

This expressiveness is consistent with the central design goal.
By shifting from a linear, event-based representation to a structural, programming-based one, the manual overhead of
repetitive event placement is removed from the source program and made explicit only in the generated MIDI artifact.
The composer works at the level of rules, patterns, and tracks, while the lowering engine materializes the temporal
event stream; the expansion ratio is simply the measurable shadow of that division of labor, and is most useful read
together with the compile-time results rather than on its own.
